% Part: model-theory
% Chapter: models-of-arithmetic
% Section: models-of-pa

\documentclass[../../../include/open-logic-section]{subfiles}

\begin{document}

\olfileid[fr]{mod}{mar}{mpa}
\section{Modèles de $\Th{PA}$}

\begin{explain}
Tout modèle non standard de~$\Th{TA}$ est également
un modèle de~$\Th{PA}$. Nous savons que de tels
modèles existent et qu'ils contiennent des
« nombres » non standards, situés au-delà de
tous les nombres standards. Comment ces éléments
sont-ils ordonnés, et combien y en a-t-il ?
La théorie plus faible~$\Th{Q}$ admet des modèles
avec un seul nombre non standard; ces !!{structure}s
simples ne sont cependant des modèles ni
de~$\Th{PA}$ ni de~$\Th{TA}$.

Pour comprendre les modèles de~$\Th{PA}$ et
de~$\Th{TA}$, examinons les faits !!{derivable}s
dans ces théories. Par exemple, $\Th{PA}$ démontre
déjà $\lforall[x][\eq/[x][x']]$ ainsi que
$\lforall[x][\lforall[y][\eq[(x+y)][(y+x)]]]$.
Ces énoncés excluent les !!{structure}s simples
où ils sont faux.

Soit $\Struct{M}$ un modèle de~$\Th{PA}$.
Chaque fois que $\Th{PA}\Proves !A$, on a
$\Sat{M}{!A}$. Notons de nouveau
$\nszero=\Assign{\Obj{0}}{M}$,
$\nssucc=\Assign{\prime}{M}$,
$\nsplus=\Assign{+}{M}$,
$\nstimes=\Assign{\times}{M}$ et
$\nsless=\Assign{<}{M}$. Un !!{sentence}~$!A$
impose alors une condition à ces interprétations.
Ainsi, si $\Sat{M}{!Q_1}$, c'est-à-dire
si
$\Sat{M}{\lforall[x][\lforall[y][
(\eq[x'][y']\lif\eq[x][y])]]}$,
la fonction $\nssucc$ doit être !!{injective}.
\end{explain}

\begin{prop}
Dans $\Struct{M}$, la relation $\nsless$ est
un ordre linéaire strict :
\begin{enumerate}
\item aucun $x\in\Domain{M}$ ne satisfait
  $x\nsless x$;
\item si $x\nsless y$ et $y\nsless z$,
  alors $x\nsless z$;
\item pour $x\neq y$, on a
  $x\nsless y$ ou $y\nsless x$.
\end{enumerate}
\end{prop}

\begin{proof}
La théorie $\Th{PA}$ démontre respectivement :
\begin{enumerate}
\item $\lforall[x][\lnot x<x]$;
\item $\lforall[x][\lforall[y][\lforall[z][
  ((x<y\land y<z)\lif x<z)]]]$;
\item $\lforall[x][\lforall[y][
  ((x<y\lor y<x)\lor\eq[x][y])]]$.\footnote{Correction éditoriale : nous supprimons la parenthèse fermante excédentaire dans le troisième énoncé affiché de la source.}
\end{enumerate}
\end{proof}

\begin{prop}
\ollabel{prop:M-discrete}
Dans l'ordre $\nsless$, $\nszero$ est le plus
petit élément de~$\Domain{M}$. Pour tout~$x$,
$x\nsless x^\nssucc$, et $x^\nssucc$ est
le plus petit élément strictement supérieur
à~$x$. Tout $x\neq\nszero$ possède un
unique $y$ tel que $y^\nssucc=x$.
On appelle $y$ le \emph{prédécesseur} de~$x$
dans~$\Struct{M}$ et on le note~$^\nssucc x$.\footnote{Correction éditoriale : la source attribue un prédécesseur à tout élément. Nous excluons $\nszero$, qui n'est le successeur d'aucun élément ; dans les blocs non standards, les prédécesseurs successifs restent non standards.}
\end{prop}

\begin{proof}
Exercice.
\end{proof}

\begin{prob}
Trouver des !!{sentence}s de~$\Lang{L_A}$
!!{derivable}s dans~$\Th{PA}$, donc vrais
dans~$\Struct{N}$, qui garantissent les
propriétés de $\nszero$, $\nssucc$ et
$\nsless$ énoncées dans la
\olref[mod][mar][mpa]{prop:M-discrete}.
\end{prob}

\begin{prop}
Tous les !!{element}s standards de~$\Struct{M}$
précèdent, selon~$\nsless$, tous ses
!!{element}s non standards.
\end{prop}

\begin{proof}
Écrivons $n$ pour
$\Value{\num{n}}{M}$, élément standard
de~$\Struct{M}$. La théorie $\Th{Q}$ démontre,
pour chaque $n\in\Nat$,
$\lforall[x][(x<\num{n}'\lif
(\eq[x][\num{0}]\lor\eq[x][\num{1}]
\lor\dots\lor\eq[x][\num{n}]))]$.
Aucun élément n'est $\nsless$-inférieur
à~$\nszero$. Si $x$ est non standard,
il ne peut donc vérifier $x\nsless n$
pour aucun~$n$. Par définition,
$x\neq n$ également. Comme $\nsless$
est un ordre linéaire, $n\nsless x$.
\end{proof}

\begin{prop}
Chaque !!{element} non standard~$x$ de
$\Domain{M}$ appartient à la suite
\[
\dots {}^{\nssucc\nssucc\nssucc}x
\nsless {}^{\nssucc\nssucc}x
\nsless {}^{\nssucc}x
\nsless x
\nsless x^{\nssucc}
\nsless x^{\nssucc\nssucc}
\nsless x^{\nssucc\nssucc\nssucc}
\nsless\dots
\]
L'ensemble ainsi obtenu est le
\emph{bloc de~$x$}, noté~$[x]$.
Il n'a ni plus petit ni plus grand élément.
Il est constitué exactement des
$y\in\Domain{M}$ pour lesquels il existe
un entier standard~$n$ tel que
$x\nsplus n=y$ ou $y\nsplus n=x$.
\end{prop}

\begin{proof}
Chaque élément non nul possède un unique
prédécesseur et chaque élément possède
un unique successeur. Les prédécesseurs
successifs d'un élément non standard
restent non standards : s'ils atteignaient
un numéral standard après un nombre fini
d'étapes, l'élément initial serait aussi
standard. On peut donc prolonger la suite
dans les deux sens sans rencontrer de borne.
Les identités
$\Th{PA}\Proves
\lforall[x][\eq[x^{\prime\dots\prime}]
[(x+\num{n})]]$,
où figurent $n$ symboles~$\prime$,
donnent la caractérisation par l'addition.
Elles montrent aussi que si $y\in[x]$,\footnote{Corrections éditoriales : la preuve distingue les prédécesseurs non nuls de zéro. Nous réparons en outre, dans la traduction, le second argument de la macro d'égalité de l'identité affichée ; la formule anglaise était correcte.}
alors $x\in[y]$ et $[x]=[y]$.
\end{proof}

\begin{prop}
Si $[x]\neq[y]$ et $x\nsless y$,
alors $u\nsless v$ pour tous
$u\in[x]$ et $v\in[y]$.
\end{prop}

\begin{proof}
Puisque les blocs sont distincts,
$x\nsplus k\nsless y$ pour tout
entier standard~$k$ : l'égalité
$x\nsplus k=y$ placerait $y$ dans~$[x]$,
et l'ordre ne peut être inversé entre
$x\nsless y$ et un de ses successeurs.
Soient $u\in[x]$ et $v\in[y]$.
Il existe des entiers standards~$m,n$
tels que $u\leq x\nsplus m$ et
$y\leq v\nsplus n$ dans l'ordre
de~$\Struct{M}$. Par conséquent,
\[
u\nsplus n
\leq x\nsplus(m+n)
\nsless y
\leq v\nsplus n.
\]
La compatibilité de l'ordre avec l'addition
et la simplification additive, démontrables
dans~$\Th{PA}$, donnent $u\nsless v$.
\end{proof}

\begin{cor}
Si $[x]\neq[y]$, alors
$[x]\cap[y]=\emptyset$.
\end{cor}

\begin{proof}
L'ordre linéaire place $x$ avant~$y$
ou $y$ avant~$x$. La proposition
précédente placerait alors tout élément
commun strictement avant lui-même,
ce qui est impossible.
\end{proof}

\begin{explain}
Les blocs distincts héritent ainsi
de l'ordre~$\nsless$ : on écrit
$[x]\nsless[y]$ lorsque $[x]\neq[y]$
et que $x\nsless y$; cela revient à
demander $u\nsless v$ pour tous
$u\in[x]$ et $v\in[y]$.
Notons $[0]$ l'ensemble des valeurs
des numéraux standards. C'est le premier
bloc; il ne rencontre aucun bloc
non standard. Deux blocs non standards
distincts ne se rencontrent pas non plus.
Un nombre fini de successeurs ou de
prédécesseurs ne permet jamais de passer
d'un bloc à un autre.\footnote{Précision éditoriale : la comparaison stricte de blocs exige $[x]\neq[y]$, condition omise dans la définition de la source. Deux éléments d'un même bloc peuvent être strictement ordonnés sans que ce bloc soit inférieur à lui-même.}
\end{explain}

\begin{prop}
Si $x$ et $y$ sont non standards,
alors $x\nsless x\nsplus y$ et
$x\nsplus y\notin[x]$.
\end{prop}

\begin{proof}
Un élément non standard~$y$ n'est pas
$\nszero$. La théorie $\Th{PA}$ démontre
$\lforall[x][\lforall[y][
(y\neq\Obj{0}\lif x<(x+y))]]$.
Si $x\nsplus y$ appartenait à~$[x]$,
l'inégalité impliquerait
$x\nsplus y=x\nsplus n^\nssucc$
pour un entier standard~$n$.
La loi de simplification additive,
démontrable dans~$\Th{PA}$, donnerait
$y=n^\nssucc$, contrairement à
l'hypothèse.
\end{proof}

\begin{prop}
Il n'existe pas de plus petit bloc
non standard.
\end{prop}

\begin{proof}
La théorie $\Th{PA}$ démontre que tout
$x$ est de la forme $y\nsplus y$
ou $(y\nsplus y)^\nssucc$ pour un~$y$.
Si $x$ est non standard, $y$ l'est aussi.
La proposition précédente donne
$y\nsless y\nsplus y$ et
$y\nsplus y\notin[y]$. Dans les deux
cas, $y\nsless x$ et $y\notin[x]$.
Le bloc~$[y]$ précède donc
strictement~$[x]$.
\end{proof}

\begin{prop}
Dans tout modèle non standard
de~$\Th{PA}$, il n'existe pas
de plus grand bloc.\footnote{Précision éditoriale : la source omet ici l'hypothèse de non-standardité. Dans le modèle standard, le seul bloc $[0]$ est bien le plus grand ; l'exercice suivant contient déjà l'hypothèse nécessaire.}
\end{prop}

\begin{proof}
Exercice.
\end{proof}

\begin{prob}
Montrer que, dans un modèle non standard
de~$\Th{PA}$, il n'existe pas de
plus grand bloc.
\end{prob}

\begin{prop}
\ollabel{prop:blocks-dense}
L'ordre des blocs est dense :
si $x\nsless y$ et $[x]\neq[y]$,
il existe un bloc~$[z]$, distinct
de ces deux blocs, situé entre eux.
\end{prop}

\begin{proof}
La théorie $\Th{PA}$ démontre que
$x\nsplus y$ est de la forme
$z\nsplus z$ ou
$(z\nsplus z)^\nssucc$.
Prenons ce $z$, la partie entière
de la moyenne de $x$ et~$y$.
Comme $[x]\neq[y]$, pour tout entier
standard~$n$, on a
$x\nsplus(2n+2)\nsless y$.
Les lois de l'ordre et de l'addition
dans~$\Th{PA}$ permettent alors
de déduire
$x\nsplus n\nsless z$ et
$z\nsplus n\nsless y$.
Ainsi $[x]\nsless[z]\nsless[y]$.\footnote{Complément éditorial : la source utilise ici un autre symbole d'addition et ne justifie pas que le point moyen sorte des deux blocs. Nous utilisons $\nsplus$ et les inégalités valables pour chaque décalage standard ; la démonstration détaillée demandée ensuite reste un exercice.}
\end{proof}

\begin{prob}
Rédiger une démonstration détaillée de la
\olref[mod][mar][mpa]{prop:blocks-dense}.
Quel !!{sentence} $\Th{PA}$ doit-elle
!!{derive} pour garantir l'existence
de~$z$ ? Pourquoi $x\nsless z$
et $z\nsless y$ ? Pourquoi enfin
$[x]\neq[z]$ et $[z]\neq[y]$ ?
\end{prob}

\begin{explain}
Dans un modèle non standard
!!{denumerable}, les blocs non standards
forment donc un ordre linéaire dense
!!{denumerable} sans extrémités.\footnote{Précision éditoriale : la source affirme cette cardinalité sans restreindre le modèle. Nous ajoutons l'hypothèse dénombrable ; chaque bloc non standard est une copie de $\Int$, et l'absence d'extrémités assure une infinité de blocs.}
Tout ordre de ce type est isomorphe
à l'ordre des rationnels. Il s'ensuit
que deux modèles non standards
!!{enumerable}s~$\Struct{M}_1$ et
$\Struct{M}_2$ de l'arithmétique vraie
ont des réduits isomorphes au langage
qui ne contient que~$<$ et l'égalité.
Les parties standards se correspondent
comme deux copies de~$\Struct{N}$.
Un isomorphisme entre les ordres
des blocs non standards s'étend
à l'intérieur de chaque bloc :
on choisit un élément dans chacun
des deux blocs correspondants, puis
on fait correspondre leurs successeurs
et leurs prédécesseurs respectifs.

Cela n'implique pas que
$\Struct{M}_1$ et $\Struct{M}_2$
soient isomorphes dans tout le langage
de l'arithmétique. L'isomorphisme
dépend du langage considéré : les
interprétations de l'addition et de la
multiplication peuvent différer.
Le théorème de Vaught, selon lequel
une théorie complète dans un langage dénombrable ne peut avoir
exactement deux modèles !!{enumerable}s
à isomorphisme près, confirme cette conclusion : l'arithmétique vraie possède
son modèle standard et un modèle non standard dénombrable, donc au moins un troisième
modèle dénombrable ; les deux modèles non standards ainsi obtenus ne sont pas isomorphes.
\end{explain}

\end{document}
